\section{A uniform polynomial-time implementation}\label{sec:runtime}

For each fixed degree threshold $k$, the preceding sections give a
finite deterministic algorithm. Its local calculations depend on $k$,
while its interval endpoints depend on the comparison history. We first
bound the size of the maintained state, then choose $k$ slowly
enough that the local calculations also take polynomial time.
Computation is measured on a fixed deterministic multitape Turing machine,
with integers represented in binary.

\subsection{Bounding the state throughout the interaction}

The first test in Algorithm~\ref{alg:compare} answers any comparison
already implied by the global comparison DAG. Such a call leaves the
stored state unchanged. Every other call starts with an incomparable pair
and makes that pair comparable permanently. Thus at most $\binom n2$
calls can change the state, however often the sorter repeats a comparison.

Lemmas~\ref{lem:preparation} and~\ref{lem:restoration} bound the number of
refinements in one call by $n+2$. Consequently the total number $P$ of
refinements, including those in a call still in progress, satisfies
\begin{equation}\label{eq:refinement-runtime}
P\le(n+2)\binom n2.
\end{equation}
Each refinement creates at most one node in the adversary's tree, so
that tree has at most $P+1$ nodes. The global comparison DAG has $n$
vertices and at most $\binom n2$ arcs. The transcript is the output stream;
its order constraints are retained in this DAG.

We store an anchor by its path in the implicit anchor tree. A path of
length $m$ gives its interval and position by $m$ halvings of $[0,1]$;
the position has $m+1$ fractional binary digits. Anchor comparisons use
these exact numerical positions.

For a wide item, Lemma~\ref{lem:median-search} bounds both the number of
search tests and the anchor precision by
\[
K_0(k)=\ceil{\log(4(B+1)R)}+2.
\]
For a narrow item, suppose the current endpoints have at most $\Lambda$
fractional digits. Its positive effective width is at least $2^{-\Lambda}$.
The selected anchor-tree interval has length $2^{-m}>W_x/8$, so
$m<\Lambda+3$. The new anchor therefore has at most $\Lambda+3$ fractional
digits. Tightening takes maxima and minima of endpoints and adds no digits.
After $P$ refinements,
\begin{equation}\label{eq:precision}
\Lambda\le K_0(k)+3P.
\end{equation}
The narrow walk makes at most $\Lambda+2$ steps, using integers with
$O(\Lambda)$ bits at each step.

\subsection{Choosing the degree threshold}\label{sec:parameter-selection}

Lemma~\ref{lem:exact-local-volume} bounds a local volume calculation by
$C_0(b)(\Lambda+1)^2$ bit operations, where the number $b$ of retained
coordinates depends only on $k$. The function $C_0$ is computable from
the finite enumeration in its proof. Finding the neighborhood and its
effective endpoints uses graph operations on at most $n$ vertices, with
a fixed polynomial cost in $n$ and $\Lambda$.

Choose a computable integer $C(k)$ that bounds $K_0(k)$ and all
parameter-dependent factors in one local calculation, including
$C_0(b)$. For powers of two $k=64,128,\ldots$, a fixed program writes the
complete parameter tuple and $C(k)$ in binary. On input size $n$, simulate
this program for $n$ steps. Among the tuples it has completed, choose the
largest $k$ satisfying
\begin{equation}\label{eq:stage-gate}
C(k)\le n.
\end{equation}
Use these parameters for the entire interaction. If no tuple qualifies,
answer all comparisons according to the fixed label order.

The initialization is polynomial: the simulation writes at most $n$
bits, and its completed tuples can be parsed and compared within a
polynomial bound. For every fixed $k$, the tuple is written in some finite
time $T(k)$. Once $n\ge\max\{T(k),C(k)\}$, it is available and affordable.
Hence the chosen $k(n)$ tends to infinity.

For a selected tuple, \eqref{eq:refinement-runtime}--\eqref{eq:stage-gate}
give $\Lambda=O(n^3)$ and a polynomial number of stored nodes. Every local
volume calculation has a fixed polynomial bit cost. A comparison call
makes at most $n+2$ refinements, and each wide refinement uses at most
$K_0(k)\le n$ estimates. Narrow refinements use the bounded walks described
above. Routing queries, propagating bounds, and testing reachability take
polynomial time in the stored state. Directed-path lengths are computed
in reverse topological order; if $H>n$, no path can require repair.
All tie rules can be implemented by scanning labels in order. Thus each
comparison is answered within one fixed polynomial number of bit
operations, independent of the chosen parameter and of the length of
the preceding interaction.

\begin{proof}[Proof of Theorem~\ref{thm:main}]
Initialization and every comparison call are deterministic and take
polynomial time. For all sufficiently large $n$, a parameter tuple
qualifies. Theorem~\ref{thm:fixed-stage} gives legality and forces
\[
q\ge\frac{n\log n-3n}{1+(2\log k(n)+18)/k(n)}
=(1-o(1))n\log n,
\]
since $k(n)\to\infty$. The fixed-order rule for the remaining input sizes
is also deterministic and legal.
\end{proof}
